% !TEX encoding = UTF-8
% Relazione aggiuntiva: Applicazione dell'inference stress test di Roccetti (2026) a RQ2 e all'Event Study
% Formato: IEEEtran (Conference Paper), coerente con bozza_paper.tex

\documentclass[conference]{IEEEtran}

\usepackage[utf8]{inputenc}
\usepackage[T1]{fontenc}
\usepackage[italian]{babel}
\usepackage{graphicx}
\usepackage{booktabs}
\usepackage{array}
\usepackage{amsmath}
\usepackage{url}
\usepackage[hidelinks]{hyperref}
\renewcommand{\arraystretch}{1.15}

\begin{document}

\title{Relazione aggiuntiva: applicazione dell'inference stress test di Roccetti (2026)\\a RQ2 e all'Event Study}

\author{
\IEEEauthorblockN{Elia Friberg}
\IEEEauthorblockA{Alma Mater Studiorum, Università di Bologna}
\and
\IEEEauthorblockN{Matteo Raggi}
\IEEEauthorblockA{Alma Mater Studiorum, Università di Bologna}
}

\maketitle

\section{Metodologia adottata}

Roccetti \cite{roccetti2026} propone un test diagnostico per conteggi rari di eventi Poisson
osservati su denominatori enormi: nei nostri dati, le nascite di un nome in un dato anno
(conteggio raro, $n$) rispetto al totale delle nascite di quell'anno/sesso/paese (denominatore
$N$, dell'ordine di 1--2 milioni per gli USA). In presenza di $N$ così grandi, l'errore standard
di Wald si restringe quasi a zero indipendentemente dalla reale solidità del segnale, producendo
$p$-value ``ultra-significativi'' che riflettono la scala del database più che l'evidenza
biologica/comportamentale sottostante.

Il rimedio proposto sostituisce l'errore standard di Wald con l'errore standard esatto di
Poisson, che dipende solo dal conteggio osservato $n$, e confronta la stima osservata con il
limite inferiore esatto dell'intervallo di confidenza di Garwood \cite{garwood1936} costruito
sullo stesso $n$. L'indice risultante è

\begin{equation}
Z_R(n) = \frac{n - \tfrac{1}{2}\chi^2(2n,\, \alpha/2)}{\sqrt{n}}
\label{eq:zr}
\end{equation}

dove $\chi^2(2n, \alpha/2)$ è il quantile $\alpha/2$ della distribuzione chi-quadrato con $2n$
gradi di libertà ($\alpha = 0{,}05$). Una proprietà algebrica notevole della formula
\eqref{eq:zr} è che il denominatore $N$ si cancella completamente: $Z_R$ dipende solo dal
conteggio grezzo $n$, non dalla popolazione totale. È proprio questa cancellazione a rendere
l'indice immune all'inflazione di significatività causata da un denominatore enorme.

\textbf{Come interpretiamo l'indice in questa relazione.} Per costruzione, $Z_R(n)$ cresce in
modo monotono e converge asintoticamente al valore normale standard $1{,}96$ (quantile normale
al $97{,}5\%$). Lo utilizziamo coerentemente con la descrizione datane dallo stesso autore
(``a standardized descriptive index of structural distance''), come \emph{indice continuo della
distanza strutturale} di ciascuna stima dal proprio riferimento asintotico $1{,}96$. Valori di
$Z_R$ più vicini a $1{,}96$ indicano conteggi la cui stima è strutturalmente più solida (meno
esposta all'illusione di precisione della scala del database); valori più lontani segnalano
conteggi la cui apparente robustezza statistica poggia su una base numerica più esigua.

\section{Applicazione a RQ2: le 17 storie di picco/crollo}

Per ciascuno dei 13 salti positivi e delle 4 flessioni verificate, abbiamo calcolato $Z_R$ sul
conteggio grezzo pre-evento e post-evento (baseline e picco), usando i conteggi reali
incrociati con \texttt{us\_names\_long.csv} / \texttt{istat\_contanomi\_full.csv} (Tabella
\ref{tab:rq2}).

\begin{table*}[t]
\centering
\footnotesize
\caption{Distanza strutturale $Z_R$ pre/post-evento per le 17 storie di picco/crollo (RQ2)}
\label{tab:rq2}
\begin{tabular}{l l r r c c}
\toprule
\textbf{nome} & \textbf{paese} & \textbf{$n_{pre}$} & \textbf{$n_{post}$} & \textbf{$Z_R$ pre} & \textbf{$Z_R$ post} \\
\midrule
Shirley  & USA     & 14.476 & 42.366 & 1,952 & 1,955 \\
Tammy    & USA     & 193    & 4.365  & 1,891 & 1,946 \\
Nakia    & USA     & 7      & 1.135  & 1,582 & 1,932 \\
Jaime    & USA     & 259    & 7.838  & 1,900 & 1,949 \\
Devante  & USA     & 10     & 131    & 1,646 & 1,876 \\
Mariah   & USA     & 423    & 1.103  & 1,914 & 1,931 \\
Nevaeh   & USA     & 8      & 1.199  & 1,607 & 1,932 \\
Jaslene  & USA     & 6      & 501    & 1,551 & 1,917 \\
Karol    & Italia  & 2      & 156    & 1,243 & 1,883 \\
Chanel   & Italia  & 8      & 63     & 1,607 & 1,838 \\
Adele    & Italia  & 449    & 1.075  & 1,915 & 1,931 \\
Elodie   & Italia  & 13     & 135    & 1,686 & 1,877 \\
Soleil   & Italia  & 98     & 474    & 1,863 & 1,916 \\
Hillary  & USA     & 2.520  & 1.064  & 1,941 & 1,931 \\
Kobe     & USA     & 1.392  & 625    & 1,934 & 1,922 \\
Alexa    & USA     & 2.002  & 708    & 1,939 & 1,924 \\
Erica    & Italia  & 863    & 416    & 1,928 & 1,913 \\
\bottomrule
\end{tabular}
\end{table*}

\textbf{Lettura.} I casi con baseline pre-evento più esiguo (Karol $n=2$, Nakia $n=7$, Jaslene
$n=6$, Chanel $n=8$, Nevaeh $n=8$, Devante $n=10$) mostrano i valori di $Z_R$ più distanti dal
riferimento asintotico sul lato ``baseline'' ($1{,}24$--$1{,}65$): sono cioè i conteggi la
cui stima di partenza è strutturalmente meno solida, coerentemente col fatto che i loro
rapporti di crescita più spettacolari (es. Nevaeh $150\times$, Jaslene $100\times$, Karol
$78\times$) nascono da un conteggio iniziale a singola/doppia cifra, in linea con l'avvertenza
già presente in tesi sul rischio di survivorship bias per i casi a bassa numerosità. I conteggi
post-evento
sono quasi tutti molto più vicini al riferimento ($1{,}84$--$1{,}96$), incluso quando il
valore assoluto resta piccolo (Chanel $n=63 \to 1{,}838$): l'evento, una volta avvenuto, tende
a essere misurato con una base numerica più solida anche quando è partito da un baseline
fragile. I quattro casi di flessione (Hillary, Kobe, Alexa, Erica) mostrano valori pre e post
entrambi ravvicinati al riferimento ($1{,}91$--$1{,}94$), partendo da conteggi già
nell'ordine delle centinaia/migliaia.

\section{Applicazione all'Event Study confermativo ($N=152$)}

Per ciascuno dei 152 eventi mediatici, $Z_R$ è stato calcolato sulla somma dei conteggi
grezzi nella finestra post-evento (fino a 2 anni), la stessa finestra usata per il calcolo
della frequenza post-evento già riportato nella tesi (Tabella \ref{tab:eventstudy}).

\begin{table*}[t]
\centering
\footnotesize
\caption{Distanza strutturale $Z_R$ (post-evento) per categoria, Event Study $N=152$}
\label{tab:eventstudy}
\begin{tabular}{l r r r r r}
\toprule
\textbf{categoria} & \textbf{$N$} & \textbf{$Z_R$ min} & \textbf{mediana} & \textbf{$Z_R$ max} & \textbf{med. $n_{post}$} \\
\midrule
Cinema \& Serie TV      & 71  & 1,455 & 1,935 & 1,956 & 1.483 \\
Musica \& Pop Culture   & 46  & 0,975 & 1,929 & 1,955 & 956   \\
Sport                   & 35  & 1,375 & 1,932 & 1,953 & 1.200 \\
\midrule
\textbf{Totale}         & \textbf{152} & \textbf{0,975} & \textbf{1,932} & \textbf{1,956} & \textbf{1.200} \\
\bottomrule
\end{tabular}
\end{table*}

Le otto storie strutturalmente più distanti dal riferimento asintotico ($Z_R$ più basso,
quelle il cui impatto netto DiD in Tabella 18 poggia sul conteggio grezzo post-evento più
esiguo) sono riportate in Tabella \ref{tab:eventstudy_fragili}.

\begin{table*}[t]
\centering
\footnotesize
\caption{Le otto storie con $Z_R$ (post-evento) più distante dal riferimento asintotico}
\label{tab:eventstudy_fragili}
\begin{tabular}{l l l r r r}
\toprule
\textbf{nome} & \textbf{categoria} & \textbf{paese} & \textbf{$n_{pre}$} & \textbf{$n_{post}$} & \textbf{$Z_R$} \\
\midrule
Celine  & Musica \& Pop Culture & IT & 168 & 1  & 0,975 \\
Lautaro & Sport                 & IT & 4   & 3  & 1,375 \\
Harry   & Cinema \& Serie TV    & IT & 3   & 4  & 1,455 \\
Nemo    & Cinema \& Serie TV    & US & 6   & 5  & 1,510 \\
Dwyane  & Sport                 & US & 10  & 8  & 1,607 \\
Andriy  & Sport                 & IT & 4   & 13 & 1,686 \\
Cher    & Musica \& Pop Culture & US & 20  & 14 & 1,696 \\
Woody   & Cinema \& Serie TV    & US & 17  & 18 & 1,728 \\
\bottomrule
\end{tabular}
\end{table*}

\textbf{Lettura.} La stragrande maggioranza dei 152 eventi (mediana $Z_R = 1{,}93$ in tutte e
tre le categorie) poggia su conteggi post-evento vicini al riferimento asintotico, cioè
strutturalmente solidi. Una minoranza di code (meno di 10 eventi su 152) mostra conteggi
post-evento a una cifra o poco più, dove la distanza dal riferimento è maggiore: un'informazione
complementare, non contraddittoria, rispetto al test di Wilcoxon già riportato nella tesi:
quest'ultimo valuta se l'effetto (variazione percentuale rispetto ai controlli) è sistematico
sull'insieme dei 152 eventi, mentre $Z_R$ valuta quanto il conteggio grezzo di ciascun singolo
evento sia strutturalmente solido come base per quella stima. Il caso ``Celine, IT, 2023''
($n_{post}=1$) è quello più esposto in assoluto.

\section{Sintesi}

Su entrambi i sottoinsiemi (RQ2 e Event Study), l'indice di distanza strutturale conferma un
quadro coerente con quanto già discusso nella tesi in termini di limiti dei casi a bassa
numerosità: i conteggi con base numerica più esigua (singola/doppia cifra) sono
sistematicamente quelli strutturalmente più distanti dal riferimento asintotico, mentre la
stragrande maggioranza dei casi (inclusi quelli con rapporti di crescita percentualmente
più eclatanti) poggia su una base numerica sufficientemente solida da collocarsi vicino al
riferimento. Questo rafforza, con una misura indipendente dal denominatore del database, un
limite già dichiarato nella tesi (il rischio di survivorship bias sui casi a bassa
numerosità), individuando puntualmente i candidati più esposti a quel rischio.

\section*{Contributi}

\begin{itemize}
  \item \textbf{Matteo Raggi:} la direzione ``da evento a nome'' dell'individuazione dei
  picchi/crolli per RQ2 (partendo dagli eventi mediatici noti per verificarne l'impatto sui
  nomi) e la relativa organizzazione per settore (Cinema \& Serie TV, Musica \& Pop Culture,
  Sport) dell'Event Study, l'intera RQ3, e la maggior parte della stesura del testo della tesi.
  \item \textbf{Elia Friberg:} RQ1, la direzione ``da nome a evento'' dell'individuazione dei
  picchi/crolli per RQ2 (partendo dalle anomalie nei dati dei nomi per risalire alla causa
  mediatica), l'applicazione dell'inference stress test di Roccetti \cite{roccetti2026}
  descritta in questa relazione (Sezioni I--IV) con relativa validazione numerica indipendente,
  e la parte restante della stesura del testo della tesi.
\end{itemize}

\begin{thebibliography}{1}

\bibitem{roccetti2026}
M.~Roccetti, ``An AI for diseases or a sick AI? The epistemological pathology of big data
overload and the loss of statistical significance,'' \emph{Frontiers in Artificial
Intelligence}, vol.~9, p.~1961525, 2026, doi: 10.3389/frai.2026.1961525 (Author's Proof).

\bibitem{garwood1936}
F.~Garwood, ``Fiducial limits for the Poisson distribution,'' \emph{Biometrika}, vol.~28,
pp.~437--442, 1936, doi: 10.1093/biomet/28.3-4.437.

\end{thebibliography}

\end{document}
